\section{System and Threat Model}\label{sec:model}

\subsection{Devices and detectors}
A fleet of IoT devices sits behind a shared network (a home, a campus wing, a factory cell). Each
device $i$ runs its own detector on its own traffic. Every packet is summarised by Kitsune's damped
incremental statistics~\cite{kitsune} (packet sizes, rates and inter-packet gaps at five time scales),
sign-log transformed, standardised with fleet-wide constants and projected onto a fleet-shared, whitened
basis of dimension $D=16$. The shared basis is what lets devices compare updates, and sixteen numbers
fit easily in one ESP-NOW frame. The detector is the online centroid detector studied by Kloft and Laskov~\cite{kloft}:
\begin{equation}
  \mathrm{score}_i(x) = \norm{x-\mu_i}/s_i, \qquad \text{alarm if } \mathrm{score}_i(x) > \tau_i,
\end{equation}
with centre $\mu_i$, spread $s_i$ and threshold $\tau_i$ set at the 99th percentile of clean
calibration scores. We write $\mu_i^0, s_i^0$ for the calibrated values and $\rho_i=\tau_i s_i^0$ for
the calibrated alarm radius, the natural unit of how far a detector can be moved.

\subsection{Self-labelled adaptation}
Learning runs in cycles. In cycle $t$, device $i$ observes a batch $X_i(t)$, treats the points it does
not alarm on as normal, and proposes an update $(d_i(t), e_i(t))$: a centre shift and a change of
log-spread. We use our signature-free rule as the local proposal: it learns only from the dense core of
accepted traffic (inside the clean 80th-percentile radius), scales a centre move by how much the outer
ring of the cloud moves with it, and allows spread growth only in proportion to outer-ring support.
Without a fleet rule, a device applies its proposal in full (\emph{single-device} adaptation). This is
the setting in which boiling-frog poisoning works~\cite{antidote,kloft}: poison that the detector accepts
moves the boundary, which lets the next poison sit further out.

\subsection{Honest drift}
Normal traffic changes. We separate two kinds. \emph{Common-mode drift} reaches many devices at once
because they share the network: a gateway or tunnel adds bytes to every packet, a route change alters
jitter, a cloud service changes its polling rate. \emph{Device-only drift} reaches one device: its own
firmware update or a change in how it is used. A detector that cannot follow drift drowns in false
alarms; one that follows everything can be poisoned.

\subsection{Threat model}\label{sec:threat}
The attacker wants an attack class (for example a Mirai scan) to stop being detected on a victim
device, or wants a defence to stop learning so honest drift turns into a flood of false alarms.
\begin{itemize}
\item \textbf{Traffic share.} The attacker injects traffic that makes up a share $p$ of the victim's
  learning batch (Kloft and Laskov's attacker share~\cite{kloft}), up to $p=0.4$. It crafts feature
  vectors freely but has no code execution on the detector and no keys.
\item \textbf{Knowledge.} White-box on the algorithm and on the victim's live state ($\mu, s$, the
  consensus), which is worst case for us.
\item \textbf{Monoculture botnet.} One exploit can compromise every device that shares a firmware and
  poison all of them in a coordinated way.
\item \textbf{Byzantine peers.} Up to $f$ fleet members are fully compromised: the attacker holds their
  keys and sends any frame it likes.
\item \textbf{Network attacker.} Without keys it can forge, replay or drop frames.
\end{itemize}
We call a victim \emph{compromised} when fewer than half of held-out real attack packets of the target
class raise an alarm. Out of scope: a compromised gateway that rewrites every device's traffic, and a
botnet that controls a majority of every peer group (we show where FOA breaks in that case).
